\section{Implementation on lattice}\label{sec:lattice}

To define a renormalized gluon quasi-PDF on the lattice, one needs to determine the renormalization factors $Z_{11}$, $Z_{22}$ and ${\overline {\delta m}}$ non-perturbatively. Each of them can be calculated separately~\cite{Zhang:2017bzy}. Alternatively, the renormalization factor can be calculated as a whole in the regularization-independent momentum subtraction scheme (RI/MOM)~\cite{Martinelli:1994ty}, as was done in the quark case in Refs.~\cite{Stewart:2017tvs,Chen:2017mzz}.

In the RI/MOM scheme, the renormalization factor of the gluon quasi-PDF can be determined by calculating the amputated Green's function of the corresponding operator in a far-offshell gluon state and requiring that the counterterm cancels all loop contributions to this Green's function at a specific gluon momentum. Schematically, one can write (a detailed investigation of the RI/MOM renormalization will be given elsewhere~\cite{Todo:2018glu})
\begin{align}
&\hspace*{-1.5em}[e^{{\overline{\delta m}}|z|}Z_{11} ^2 Z_3]^{-1}(z n,p^R_z, 1/a, \mu_R)=\non\\
&\hspace*{-2em}\left.\frac{P_{ij}^{ab} \langle 0| T [A^{a,i} (p)  O^{1}(z,0)A^{b,j}(-p) ]|0\rangle|_{amp}}{P_{ij}^{ab} \langle 0|T [A^{a,i} (p)   O^{1}(z,0)A^{b,j} (-p)]|0\rangle_{amp,tree}}\right|_{\tiny\begin{matrix}p^2=-\mu_R^2 \\ \!\!\!\!p_z=p^R_z\end{matrix}},
\end{align}
in the case of $O^1$. Here $A^{a,i}(p)$ denotes the momentum space gluon field with momentum $p$. $Z_3$ is the gluon field renormalization constant. For $O^{2,3,4}$, $Z_{11} ^2$ should be replaced by $Z_{22} ^2$, $Z_{11}Z_{12}$, and $Z_{22} ^2$, respectively. $P_{ij}^{ab}$ is a projection operator with color indices $a,b$ and Lorentz indices $i, j$. A simple example of the projection is $P_{ij}^{ab}= \delta^{ab}g_{\perp,ij}$.
$\mu_R$ and $p_z^R$ are unphysical scales introduced in the RI/MOM scheme~\cite{Stewart:2017tvs,Chen:2017mzz}.

After the non-perturbative renormalization, we can write down the factorization for the renormalized gluon/quark quasi-PDFs:
\begin{align}\label{allfac}
&\tilde g(x ,P_z, p_z^R, \mu_R)=\int_0^1\frac{dy}{y}\Big[C_{gg}\Big(\frac{x}{y}, \frac{\mu_R}{p_z^R},\frac{y P_z}{\mu}, \frac{y P_z}{p_z^R}\Big)g(y, \mu)\non\\
&+C_{gq}\Big(\frac{x}{y}, \frac{\mu_R}{p_z^R},\frac{y P_z}{\mu}, \frac{y P_z}{p_z^R}\Big)q_j(y, \mu)\Big]+\mathcal O\Big(\frac{M^2}{P_z^2}, \frac{\Lambda_{\rm QCD}^2}{P_z^2}\Big),\nonumber\\
&\tilde q_i(x ,P_z, p_z^R, \mu_R)=\int_0^1\frac{dy}{y}\Big[C_{q_i q_j}\Big(\frac{x}{y}, \frac{\mu_R}{p_z^R},\frac{y P_z}{\mu}, \frac{y P_z}{p_z^R}\Big)q_j(y, \mu)\non\\
&+C_{qg}\Big(\frac{x}{y}, \frac{\mu_R}{p_z^R},\frac{y P_z}{\mu}, \frac{y P_z}{p_z^R}\Big)g(y, \mu)\Big]+\mathcal O\Big(\frac{M^2,}{P_z^2}, \frac{\Lambda_{\rm QCD}^2}{P_z^2}\Big),
\end{align}
where the momentum fraction of all parton distributions is within $[0,1]$. A summation of $j$ over all quark/antiquark species is implied. Like the renormalized quasi-PDFs, the hard coefficients $C_{gg}$, $C_{qg}$, $C_{q_i q_j}$ and $C_{gq}$ also depend on unphysical scales, e.g. $\mu_R, p_z^R$ and the hadron momentum $P_z$. But PDFs extracted from Eq.~(\ref{allfac}) are independent of these scales. For polarized quasi-PDFs, the factorization takes a similar form as Eq.~(\ref{allfac}).
Details of the proof of the above factorization formula using the operator product expansion~\cite{Izubuchi:2018srq} as well as the perturbative calculation of the hard coefficients at one-loop accuracy will be presented in a separate publication~\cite{Todo:2018glu}.
